\section{Limitations}
\label{sec:limitations}

% Provenance wording (generated by make_tables.py); fallbacks for standalone compile.
\providecommand{\ArtifactPosCtrlMode}{the fast-exit positive control is an in-process \texttt{shadow} rehearsal}
\providecommand{\ArtifactPairedMode}{the paired race is an in-process \texttt{shadow} rehearsal}
\providecommand{\ArtifactRqEightRemains}{to replace the in-process rehearsal of the paired race with the real-client-code localnet run}

We are explicit about what this work does and does not establish.

\emph{(1) Conditional, not universal.} The attack requires the joint conditions of
Table~\ref{tab:applicability}: a small revealed committee, an open settlement window
($\Tsettle<\Tacc$), a reachable value sink, and positive bribe-adjusted profit. Designs that enforce
$\Tsettle\ge\Tacc$ are immune by construction (Propositions~\ref{prop:immunity},~\ref{prop:reconf-safe}).

\emph{(2) Necessary, not sufficient.} The $F{+}1$ equivocators are the quorum-intersection floor;
completion additionally requires an honest-vote split, whose probability $\psucc$ depends on
protocol/network scheduling and, at the floor, is modest ($\ArtifactPsuccFloor$ at $N=31$) unless the
adversary controls views. Our profitability analyses fix $\psucc=0.85$, which assumes that
view-control regime (or a $1.45\times$-larger bribe, $k\approx F{+}6$, to reach it by purchase); we
bracket those results down to the floor rather than present $0.85$ as expected
(Figure~\ref{fig:psucc}, \S\ref{sec:floor}).

\emph{(3) $q$ is load-bearing; we ground it but do not measure it on a live network.} The bribe price
is a fraction $q$ of stake-at-risk, and $q$ depends on accountability latency, unbonding/withdrawal
delay, evidence validity, and enforcement reliability. We sweep $q$, measure the race in a
discrete-event testbed, and evaluate $q$ at the documented parameters of six deployed systems
(\S\ref{sec:eval}, where the documented withdrawal delay gives $q\approx1$). To close this, \S\ref{sec:measured}
specifies and \emph{pre-registers} a measurement protocol (\texttt{prereg/rq8-v1}) that obtains
$q$, $\Tacc$, and $\Tsettle$---with cooperative and adversarial bounds on $\Tacc$---on a real safe
path (Cosmos/CometBFT with finality-gated IBC) and a real open path (a fast intent bridge), each
without executing the attack on any shared network, together with a frozen falsification rule and a
per-system distance-to-cliff. We have executed the \emph{passive Cosmos arm} and the \emph{bridge
settlement marginal} on real data---grounding the accountability side on real double-sign incidents
(the operative deadline is the evidence-expiry horizon, not unbonding, and it is soft: enforcement
reached ${\sim}25$ days, our adversarial $\Tacc^{\mathrm{hi}}$) and exhibiting a real seconds-scale
open window ($\Tsettle\!\ll\!\Tacc$) on a live value path---and we demonstrate the paired race, the
fast-exit positive control, and the dual verdict end-to-end on the validated testbed
(\ArtifactPairedMode, and \ArtifactPosCtrlMode), with a per-system distance-to-cliff. We have
additionally measured the native-IBC negative control on a real public IBC channel
(Cosmos\,Hub/Osmosis), which returns \textsc{ruled-out}. What remains is \ArtifactRqEightRemains, and to
re-derive the Cosmos figures against a self-run archive node; $q_{\max}$ stays an enforcement-conditional
\emph{upper bound} throughout. Completing that is the remaining empirical step.

\emph{(4) Recruitment and coordination.} The model assumes rational validators
(Assumption~\ref{ass:rational}) and abstracts the problem of recruiting $F{+}1$ specific committee
members within an epoch; collusion-resistance and whistleblowing dynamics
(cf.~\cite{kelkar2025omerta}) may materially change costs.

\emph{(5) Strong accountability defeats it.} Long unbonding, valid and fast evidence, persistent
slashability across rotation, and accountability-gated cross-shard acceptance all reduce profitability;
where they jointly hold, the honest framing of this work is a design rule rather than a live threat.

\emph{(6) Reconfiguration cuts both ways.} Reconfiguration may widen or close the window depending on
state/receipt handoff (\S\ref{sec:reconf}); we characterize both regimes but do not claim rotation
generically helps the attacker.

\emph{(7) Local model and testbed only.} The evaluation is an analytical model, Monte-Carlo
simulation, a discrete-event testbed with real Ed25519 signatures, and a case study on documented
public parameters. It targets no live system, and the real-systems figures are approximate public
values to be confirmed against current specifications before camera-ready.
